\documentclass[10pt,a4paper]{article}
\usepackage[utf8]{inputenc}
\usepackage[T1]{fontenc}
\usepackage[french]{babel}
\usepackage[margin=1.5cm]{geometry}
\usepackage{lmodern}
\usepackage{amsmath}
\usepackage{multicol}

% --- Définitions ---
\renewcommand{\familydefault}{\sfdefault}
\setlength{\columnsep}{1cm}
\pagestyle{empty}
\newcommand{\code}[1]{\texttt{#1}}

\begin{document}

\begin{center}
    \huge\textbf{Fiche Récapitulative d'Algorithmique}
\end{center}
\vspace{0.4cm}

\begin{multicols}{2}

% --- STRUCTURE ---
\section*{Structure Générale}
Un algorithme est une suite d'opérations visant à résoudre un problème. Le pseudo-code le décrit en langage simple.
\hrule
\vspace{1mm}
\textbf{ALGORITHME} Nom\_Algorithme \\
\textbf{Variables} \\
\quad \textit{...déclarations...} \\
\textbf{DEBUT} \\
\quad \textit{...instructions...} \\
\textbf{FIN}

% --- ALGORIGRAMMES ---
\section*{Algorigrammes (Symboles)}
Représentation graphique d'un algorithme.
\hrule
\vspace{1mm}
\begin{tabular}{ll}
    \textbf{Forme} & \textbf{Signification} \\
    \hline
    Ovale & Début ou Fin \\
    Rectangle & Traitement, Action \\
    Losange & Condition, Test \\
    Parallélogramme & Entrée ou Sortie (Lire/Ecrire) \\
    \hline
\end{tabular}

% --- VARIABLES ---
\section*{Variables et Types}
Une variable stocke une donnée. Elle est définie par un \textbf{nom} et un \textbf{type}.
\hrule
\begin{itemize}
    \item \textbf{Numérique} : Nombres entiers ou réels. \\ \textit{Ex: Variable age en Numérique}
    \item \textbf{Chaîne} : Texte (entre guillemets). \\ \textit{Ex: Variable nom en Chaîne}
    \item \textbf{Booléen} : VRAI ou FAUX. \\ \textit{Ex: Variable estMajeur en Booléen}
\end{itemize}

% --- OPERATIONS ---
\section*{Opérations de Base}
\hrule
\begin{description}
    \item[Affectation] Stoque une valeur. \\ \code{variable $\leftarrow$ valeur} \textit{ ou } \code{variable = valeur}
    \item[Lecture] Récupère une entrée utilisateur. \\ \code{Lire variable}
    \item[Écriture] Affiche une valeur ou un message. \\ \code{Ecrire "Bonjour"}, \code{Ecrire variable}
\end{description}

\subsection*{Opérateurs}
\begin{itemize}
    \item \textbf{Arithmétiques} : +, -, *, /, \textasciicircum{} (puissance)
    \item \textbf{Comparaison} : =, <>, <, >, <=, >= 
    \item \textbf{Logiques} : ET, OU, NON
\end{itemize}

\columnbreak

% --- STRUCTURES DE CONTROLE ---
\section*{Structures de Contrôle}

\subsection*{1. La Conditionnelle ("Si")}
Exécute un code si une condition est vraie.
\hrule
\vspace{1mm}
\textbf{Si} <condition> \textbf{Alors} \\
\quad \textit{... instructions si VRAI ...} \\
\textbf{Sinon} \\
\quad \textit{... instructions si FAUX ...} \\
\textbf{Fin Si}

\subsection*{2. La Boucle "Tant Que"}
Répète un bloc d'instructions \textbf{tant que} la condition est VRAIE. La condition est testée au début.
\hrule
\vspace{1mm}
\textbf{Tant que} <condition> \textbf{faire} \\
\quad \textit{... instructions à répéter ...} \\
\textbf{Fin TantQue}

\subsection*{3. La Boucle "Pour"}
Répète un bloc pour un nombre \textbf{connu} d'itérations.
\hrule
\vspace{1mm}
\textbf{Pour} compteur de \textit{début} à \textit{fin} \textbf{faire} \\
\quad \textit{... instructions à répéter ...} \\
\textbf{Fin Pour}

\subsection*{4. La Boucle "Répéter"}
Répète un bloc \textbf{jusqu'à} ce que la condition devienne VRAIE. S'exécute au moins une fois.
\hrule
\vspace{1mm}
\textbf{Répéter} \\
\quad \textit{... instructions à répéter ...} \\
\textbf{Jusqu`a} <condition>

% --- TABLEAUX ---
\section*{Les Tableaux}
Collection de valeurs de \textbf{même type}, accessibles via un \textbf{indice}.
\hrule
\begin{description}
    \item[Déclaration] \textbf{Tableau} nomTableau[taille] \textbf{en} Type \\ \textit{Ex: Tableau notes[19] en Numérique (20 notes, de 0 à 19)}
    \item[Accès] Se fait via l'indice (position). \\ \code{notes[0] $\leftarrow$ 15} \textit{ (affecte 15 à la première case)}
    \item[Parcours] Souvent avec une boucle \textbf{Pour}. \\
    \textbf{Pour} i de 0 à 19 \textbf{faire} \\
    \quad Ecrire notes[i] \\
    \textbf{Fin Pour}
\end{description}

\end{multicols}
\end{document}
